\section{Introduction}
\label{sec:intro}

Accurate presurgical diagnosis of brain tumors is essential for treatment planning and patient management. Different tumor types require substantially different therapeutic strategies---gliomas may necessitate maximal safe surgical resection, meningiomas are often managed with surgery or observation, while brain metastases may be treated with radiation therapy or targeted systemic treatment~\cite{price2024}. A reliable, automated preoperative classification system can thus play an important role in supporting clinical decision-making and improving patient outcomes.

Magnetic Resonance Imaging (MRI) is the primary imaging modality used for brain tumor diagnosis and monitoring~\cite{wang2024}. In clinical practice, radiologists interpret multi-sequence MRI findings together with patient demographics and summarize their assessment in radiology reports. However, manual interpretation is challenged by overlapping imaging characteristics across tumor types, inter-observer variability, and the complexity of integrating information from multiple modalities. This motivates the development of machine learning models that can systematically combine multimodal data for tumor classification.

In this project, we address the task of five-class brain tumor classification using multimodal data comprising radiology reports, MRI-derived features, and clinical information. The classification task presents several fundamental challenges:

\begin{itemize}
  \item \textbf{Severe class imbalance:} Gliomas account for nearly half of all cases (46.6\%), while Pineal/Choroid Plexus tumors constitute only 1.1\%, creating a 40:1 imbalance ratio that biases models toward majority classes.
  \item \textbf{Multimodal integration:} Diagnostic information is spread across modalities with fundamentally different statistical properties---unstructured text (radiology reports), high-dimensional continuous features (image embeddings), categorical variables (signal intensity patterns), and tabular data (demographics)---requiring careful fusion strategies.
  \item \textbf{Distribution shift:} The test set may exhibit systematic differences from training data in feature distributions, missing data patterns, and location vocabulary. Standard cross-validation estimates may be overly optimistic under such covariate shift, motivating the use of robustness-promoting strategies throughout our design.
  \item \textbf{Limited sample size:} With only 2,266 training samples and feature dimensions up to 8,192 (raw image features), there is a severe risk of overfitting that demands aggressive dimensionality reduction and regularization.
\end{itemize}

Our approach is guided by a principled analysis of what each data modality contributes. Before designing any model, we first measure the solo discriminative power of each source through 5-fold cross-validation: radiology reports ($\sim$84\% Micro-F1), clinical features ($\sim$67\%), image features ($\sim$57\%), and raw radiomic features ($\sim$52\%, near the 46.6\% majority-class baseline). This ranking directly informs our architecture---we allocate the most capacity to the strongest signal (reports) and remove modalities that add noise (raw radiomics, shifted image features in the neural network).

Based on this analysis, we design a multimodal ensemble system with the following key components:

\begin{enumerate}
  \item \textbf{Dual-path neural network with cross-modal gating.} Path A encodes radiology reports using BioClinicalBERT~\cite{alsentzer2019} (fine-tuning the top 2 layers). Path B encodes tabular features (clinical data, signal intensity patterns, tumor location with hierarchical parsing). A learned gate controls how much cross-modal information to inject, allowing the model to rely on reports when they are informative and supplement with tabular data when reports are ambiguous.
  \item \textbf{Tree-based models on report and clinical features.} Three gradient-boosted tree models (XGBoost~\cite{chen2016}, LightGBM~\cite{ke2017}, CatBoost~\cite{prokhorenkova2018}) are trained on TF-IDF vectorized report text combined with clinical and location features. These models capture different decision boundary patterns from the neural network, providing ensemble diversity.
  \item \textbf{Weighted probability ensemble.} A grid search over OOF predictions determines optimal blend weights across all 6 models, finding NN=40\% and MLU=60\% as optimal. The final ensemble achieves \best{87.25\%} out-of-fold Micro-F1 with grid-search weights, and 86.98\% with the submission configuration.
\end{enumerate}

The remainder of this report is organized as follows. Section~\ref{sec:dataset} describes the dataset and exploratory data analysis. Section~\ref{sec:method} presents our methodology, including feature engineering, model architecture design, and ensemble strategy, with emphasis on \emph{why} each design choice was made. Section~\ref{sec:experiments} details experiments and results with ablation studies. Section~\ref{sec:discussion} discusses key findings, challenges, and lessons learned. Section~\ref{sec:conclusion} concludes with future directions.
